\documentclass{article}
\usepackage{PRIMEarxiv} % Core package for the PrimeAI Template
\usepackage{amsmath, amssymb, amsthm, bm, dcolumn} % Add amsthm here for the proof environment
\usepackage[numbers,sort&compress]{natbib} % Natbib for citations
\usepackage{graphicx} % For high-quality images
\usepackage{multicol} % Multi-column figures/tables
\usepackage{hyperref} % Hyperlinks
\usepackage{listings} % Code listings
\usepackage{authblk} % For structured affiliations
% Define theorem style (optional)
\theoremstyle{plain}
\newtheorem{theorem}{Theorem}
\renewcommand{\qedsymbol}{}

% Custom commands for primes and qubit encoding
\newcommand{\primeQubit}[1]{|p_{#1}\rangle}
\newcommand{\primeGate}[1]{U_{p_{#1}}}

\usepackage{wrapfig}
\usepackage[pscoord]{eso-pic}
\usepackage[fulladjust]{marginnote}
\reversemarginpar

% Typesetting improvements without footnote patching
\usepackage[protrusion=true, expansion=true, tracking=false]{microtype}
\microtypecontext{spacing=nonfrench}

% Line numbers
\usepackage[right]{lineno}

% Text layout - adjust as needed
\raggedright
\setlength{\parindent}{0.5cm}
\textwidth 5.25in 
\textheight 8.75in

% Set double spacing
\usepackage{setspace} 
\doublespacing

% Adjust width for specific content
\usepackage{changepage}

% Adjust caption style
\usepackage[aboveskip=1pt,labelfont=bf,labelsep=period,singlelinecheck=off]{caption}

% Remove brackets from references
\makeatletter
\renewcommand{\@biblabel}[1]{\quad#1.}
\makeatother

% Header, footer, and page numbers
\usepackage{lastpage,fancyhdr}
\pagestyle{fancy}
\fancyhf{}
\fancyhead[L]{Preprint - PrimeAI Enhanced Template}
\fancyfoot[C]{\scriptsize Multiplicity Theory © 2024 Ryan Van Gelder - Citizen Gardens \\ Licensed Under MIT and CC BY-NC-SA 4.0.}
\fancyfoot[R]{Page \thepage\ of \pageref{LastPage}}
\renewcommand{\footrule}{\hrule height 2pt \vspace{2mm}}

\title{Mathematical Framework Integrating Prime Regularities with Multiplicity Theory}
\author[1]{Miroslav Zidek}
\author[2]{Ryan Van Gelder}
\affil{Citizen Gardens - The Foundation of Multiplicity\\info@citizengardens.org}
\date{2025}

\begin{document}
\maketitle

\begin{abstract}
This paper presents a cohesive mathematical framework integrating the inherent properties of prime numbers with the dynamic principles of Multiplicity Theory. The aim is to unify the discrete regularities of primes with the emergent, interconnected behaviors modeled in Multiplicity. By leveraging tools such as eigenvalues, tensor networks, and recursive feedback, this framework provides a robust foundation for exploring complex systems across quantum computing, cryptography, and interdisciplinary domains.
\end{abstract}
\section{Introduction}
Prime numbers, as the fundamental building blocks of number theory, exhibit unique distributional regularities. Multiplicity Theory, on the other hand, emphasizes the interconnectedness and emergent dynamics of systems across scales. This framework synthesizes these perspectives, leveraging primes as discrete states and embedding their interactions within a dynamic multiplicative structure.

\section{Key Mathematical Constructs}

\subsection{Prime-Based Encoding}
Primes serve as unique identifiers for states within the Multiplicity framework. Each prime $p_i$ represents an eigenvalue, encoding stability or fundamental modes of a system:
\begin{equation}
\phi(p_i, t) = p_i \cdot F(t),
\end{equation}
where $F(t)$ is a feedback function modulating the prime's influence over time.

\subsection{Dynamic Multiplicity Equation}
The enhanced Multiplicity equation incorporates time-dependent and recursive interactions:
\begin{equation}
\frac{\partial \rho_k}{\partial t} = \alpha_k(t) \rho_k + \beta_k(t) I_k + \gamma_k(t) \sum_{j} T_{kj} \rho_j + \lambda(t) \left(\Omega_B(\rho) + \Omega_{FS}(\rho)\right) + \eta_k \rho_k^2 + \xi_k(t).
\end{equation}
Here:
\begin{itemize}
    \item $\alpha_k, \beta_k, \gamma_k$: Time-dependent interaction coefficients.
    \item $T_{kj}$: Tensor representing interaction between modes.
    \item $\Omega_B$ and $\Omega_{FS}$: Geometric terms dependent on state $\rho$.
    \item $\xi_k(t)$: Stochastic term capturing noise.
\end{itemize}

\subsection{Recursive Feedback and Prime Dynamics}
Recursive feedback captures evolving relationships among primes:
\begin{equation}
M(t) = \sum_{i=1}^N \left( p_i \cdot \cos(\omega_i t + \phi_i) \cdot F(t) \cdot \left(1 + \xi_i(t)\right) \right) + \prod_{j=1}^M \left(p_j(t) + \epsilon_j\right),
\end{equation}
where:
\begin{itemize}
    \item $\omega_i$: Angular frequency of prime interactions.
    \item $\phi_i$: Phase shift associated with $p_i$.
    \item $\xi_i(t)$: Time-dependent noise affecting $p_i$.
    \item $\epsilon_j$: Stochastic variability term for prime product interactions.
\end{itemize}

\subsection{Tensor Network Representation}
Relationships between primes can be captured using tensor networks:
\begin{equation}
T_{ijk} = \sum_{a,b,c} \phi_a \phi_b \phi_c \quad \text{where } \phi_a = e^{2\pi i n_a/p_a}.
\end{equation}
This formulation captures multidimensional dependencies among primes and their encoded states.

\subsection{Quantum State Integration}
Incorporating primes into quantum systems, the state evolution is modeled as:
\begin{equation}
|\psi(t)\rangle = \sum_{i=1}^N \alpha_i(t) |p_i\rangle + \epsilon(t),
\end{equation}
where:
\begin{itemize}
    \item $\alpha_i(t)$: Amplitude for state $|p_i\rangle$.
    \item $\epsilon(t)$: Stochastic noise term.
\end{itemize}

\section{Applications and Implications}

\subsection{Cryptography}
Prime-based encoding enhances quantum-resistant cryptographic schemes:
\begin{equation}
E(x,t) = H\left(\prod_{i=1}^N p_i^{x_i}\right) \cdot F(t),
\end{equation}
where $H$ represents a secure hashing function and $F(t)$ provides time-dependent modulation.

\subsection{Quantum Computing}
The prime-eigenvalue framework supports scalable quantum algorithms by embedding primes into eigenvalue dynamics. Tensor networks allow efficient simulation of high-dimensional systems.

\subsection{Interdisciplinary Insights}
This framework bridges domains such as social physics and neuroscience by modeling emergent behaviors from fundamental units, akin to prime distributions in number theory.

\section{Conclusion}
By unifying prime regularities with the interconnected principles of Multiplicity Theory, this framework offers a novel perspective on complex systems. Future work will focus on experimental validation and interdisciplinary applications.


\nocite{*}
\bibliographystyle{unsrt}
\bibliography{references}

\end{document}